% GENERATED FILE. Edit canonical lessons, then run npm run generate:books.
% canonical-source-hash: fnv1a64:06af3d84bb017bdb
% canonical-lessons: MW-C285-akhar, MW-C285-pano, MW-C285-path, MW-C285-ratno, MW-C285-ucharno

\chapter{Learning by Heart}
\label{ch:mw-learning-by-heart}
\hlchaptermodality{285}

\begin{chapteropening}
\textbf{By the end of this chapter:} \emph{I can say how I learn a lesson.}

Complete the last lesson of chapter 285: I can say how I learn a lesson.
\end{chapteropening}

\section[\texorpdfstring{\mw{आखर} (ākhar)}{आखर (ākhar)}]{\mw{आखर} (ākhar) — a letter of the alphabet}

\label{lesson:MW-C285-akhar}



\begin{quote}
Before the new one: say the Marwadi for an inkpot, then the Marwadi for to study.
\end{quote}

\subsection*{You'll want to know: \mw{आखर}}
\textbf{\mw{आखर}} — \emph{ākhar} — \textquotedblleft{}a letter of the alphabet\textquotedblright{}.

\begin{quote}
\mw{यो} \mw{आखर} \mw{नुवो} \mw{है।}
\end{quote}

\begin{quote}
\emph{yo ākhar nuvo hai.}
\end{quote}

\begin{quote}
This letter is new.
\end{quote}

A masculine noun.

\begin{cousinweb}
\textbf{\mw{आखर}} (\emph{ākhar}) is inherited from Sanskrit \emph{akṣara}, \textquotedblleft{}a syllable, a letter\textquotedblright{}; CDIAL gives no Marwari form in this entry, but lists Hindi \emph{ākhar}, \textquotedblleft{}letter\textquotedblright{} (CDIAL 38).
\end{cousinweb}

\begin{grammarlens}[title={What you've built}]
One more word for this chapter.
\end{grammarlens}

\subsection*{Your turn}
\begin{itemize}
\raggedright
  \item \emph{Say it:} \emph{ākhar}
  \item \emph{Say it:} \emph{ākhar}, once more
  \item \emph{Say it:} \emph{yo ākhar nuvo hai.} — This letter is new.
  \item \emph{Recall:} say the Marwadi for an inkpot, then the Marwadi for to study, then say \emph{ākhar} again
\end{itemize}

\begin{tcolorbox}[breakable,colback=teal!4,colframe=teal!35!black,title={Before you move on}]
What does \mw{आखर} mean? (\textquotedblleft{}A letter of the alphabet\textquotedblright{}.) Say it once more.
\end{tcolorbox}
\section[\texorpdfstring{\mw{पानो} (pāno)}{पानो (pāno)}]{\mw{पानो} (pāno) — a page}

\label{lesson:MW-C285-pano}



\begin{quote}
Before the new one: say the Marwadi for to study, then the Marwadi for a letter of the alphabet.
\end{quote}

\subsection*{You'll want to know: \mw{पानो}}
\textbf{\mw{पानो}} — \emph{pāno} — \textquotedblleft{}a page\textquotedblright{}.

\begin{quote}
\mw{किताब} \mw{मांय} \mw{सो} \mw{पाना} \mw{है।}
\end{quote}

\begin{quote}
\emph{kitāb mā̃y so pānā hai.}
\end{quote}

\begin{quote}
The book has a hundred pages.
\end{quote}

A masculine noun.

\begin{cousinweb}
\textbf{\mw{पानो}} (\emph{pāno}) is inherited from Sanskrit \emph{parṇa}, \textquotedblleft{}a leaf\textquotedblright{}; CDIAL lists Marwari \emph{pān}, \textquotedblleft{}leaf\textquotedblright{}, Sindhi \emph{pano}, \textquotedblleft{}page of book\textquotedblright{}, and Hindi \emph{pannā}, \textquotedblleft{}leaf, page\textquotedblright{} (CDIAL 7918).
\end{cousinweb}

\begin{grammarlens}[title={What you've built}]
One more word for this chapter.
\end{grammarlens}

\subsection*{Your turn}
\begin{itemize}
\raggedright
  \item \emph{Say it:} \emph{pāno}
  \item \emph{Say it:} \emph{pāno}, once more
  \item \emph{Say it:} \emph{kitāb mā̃y so pānā hai.} — The book has a hundred pages.
  \item \emph{Recall:} say the Marwadi for to study, then the Marwadi for a letter of the alphabet, then say \emph{pāno} again
\end{itemize}

\begin{tcolorbox}[breakable,colback=teal!4,colframe=teal!35!black,title={Before you move on}]
What does \mw{पानो} mean? (\textquotedblleft{}A page\textquotedblright{}.) Say it once more.
\end{tcolorbox}
\section[\texorpdfstring{\mw{पाठ} (pāṭh)}{पाठ (pāṭh)}]{\mw{पाठ} (pāṭh) — a lesson}

\label{lesson:MW-C285-path}



\begin{quote}
Before the new one: say the Marwadi for a letter of the alphabet, then the Marwadi for a page.
\end{quote}

\subsection*{You'll want to know: \mw{पाठ}}
\textbf{\mw{पाठ}} — \emph{pāṭh} — \textquotedblleft{}a lesson\textquotedblright{}.

\begin{quote}
\mw{यो} \mw{पाठ} \mw{सोरो} \mw{है।}
\end{quote}

\begin{quote}
\emph{yo pāṭh soro hai.}
\end{quote}

\begin{quote}
This lesson is easy.
\end{quote}

A masculine noun.

\begin{cousinweb}
\textbf{\mw{पाठ}} (\emph{pāṭh}) is a tatsama, taken from Sanskrit \emph{pāṭha}, \textquotedblleft{}recitation, study\textquotedblright{} (Monier-Williams).
\end{cousinweb}

\begin{grammarlens}[title={What you've built}]
One more word for this chapter.
\end{grammarlens}

\subsection*{Your turn}
\begin{itemize}
\raggedright
  \item \emph{Say it:} \emph{pāṭh}
  \item \emph{Say it:} \emph{pāṭh}, once more
  \item \emph{Say it:} \emph{yo pāṭh soro hai.} — This lesson is easy.
  \item \emph{Recall:} say the Marwadi for a letter of the alphabet, then the Marwadi for a page, then say \emph{pāṭh} again
\end{itemize}

\begin{tcolorbox}[breakable,colback=teal!4,colframe=teal!35!black,title={Before you move on}]
What does \mw{पाठ} mean? (\textquotedblleft{}A lesson\textquotedblright{}.) Say it once more.
\end{tcolorbox}
\section[\texorpdfstring{\mw{रटणो} (raṭṇo)}{रटणो (raṭṇo)}]{\mw{रटणो} (raṭṇo) — to learn by heart}

\label{lesson:MW-C285-ratno}



\begin{quote}
Before the new one: say the Marwadi for a page, then the Marwadi for a lesson.
\end{quote}

\subsection*{You'll want to know: \mw{रटणो}}
\textbf{\mw{रटणो}} — \emph{raṭṇo} — \textquotedblleft{}to learn by heart\textquotedblright{}.

\begin{quote}
\mw{छोरो} \mw{पाठ} \mw{रटै} \mw{है।}
\end{quote}

\begin{quote}
\emph{chhoro pāṭh raṭai hai.}
\end{quote}

\begin{quote}
The boy learns the lesson by heart.
\end{quote}

Past (I …): \mw{म्हैं} \mw{पाठ} \mw{रट्यो}, \emph{mhaĩ pāṭh raṭyo}.

\begin{cousinweb}
\textbf{\mw{रटणो}} (\emph{raṭṇo}) is inherited from Sanskrit \emph{raṭati}, \textquotedblleft{}cries, yells\textquotedblright{}; CDIAL gives no Marwari form in this entry, but lists Hindi \emph{raṭnā}, \textquotedblleft{}to cry out, repeat\textquotedblright{} (CDIAL 10590). Saying a thing over and over is how it is learnt by heart.
\end{cousinweb}

\begin{grammarlens}[title={What you've built}]
One more word for this chapter.
\end{grammarlens}

\subsection*{Your turn}
\begin{itemize}
\raggedright
  \item \emph{Say it:} \emph{raṭṇo}
  \item \emph{Say it:} \emph{raṭṇo}, once more
  \item \emph{Say it:} \emph{chhoro pāṭh raṭai hai.} — The boy learns the lesson by heart.
  \item \emph{Recall:} say the Marwadi for a page, then the Marwadi for a lesson, then say \emph{raṭṇo} again
\end{itemize}

\begin{tcolorbox}[breakable,colback=teal!4,colframe=teal!35!black,title={Before you move on}]
What does \mw{रटणो} mean? (\textquotedblleft{}To learn by heart\textquotedblright{}.) Say it once more.
\end{tcolorbox}
\section[\texorpdfstring{\mw{उचारणो} (uchārṇo)}{उचारणो (uchārṇo)}]{\mw{उचारणो} (uchārṇo) — to pronounce}

\label{lesson:MW-C285-ucharno}



\begin{quote}
Before the new one: say the Marwadi for a lesson, then the Marwadi for to learn by heart.
\end{quote}

\subsection*{You'll want to know: \mw{उचारणो}}
\textbf{\mw{उचारणो}} — \emph{uchārṇo} — \textquotedblleft{}to pronounce\textquotedblright{}.

\begin{quote}
\mw{छोरो} \mw{आखर} \mw{उचारै} \mw{है।}
\end{quote}

\begin{quote}
\emph{chhoro ākhar uchārai hai.}
\end{quote}

\begin{quote}
The boy pronounces the letters.
\end{quote}

Past (I …): \mw{म्हैं} \mw{सबद} \mw{उचार्यो}, \emph{mhaĩ sabad uchāryo}.

\begin{cousinweb}
\textbf{\mw{उचारणो}} (\emph{uchārṇo}) is inherited from Sanskrit \emph{uccārayati}, \textquotedblleft{}utters\textquotedblright{}; CDIAL lists Old Marwari \emph{ūcaraï}, \textquotedblleft{}recites\textquotedblright{}, and Hindi \emph{ucārnā}, \textquotedblleft{}to pronounce\textquotedblright{} (CDIAL 1641).
\end{cousinweb}

\begin{grammarlens}[title={What you've built}]
One more word for this chapter.
\end{grammarlens}

\subsection*{Your turn}
\begin{itemize}
\raggedright
  \item \emph{Say it:} \emph{uchārṇo}
  \item \emph{Say it:} \emph{uchārṇo}, once more
  \item \emph{Say it:} \emph{chhoro ākhar uchārai hai.} — The boy pronounces the letters.
  \item \emph{Recall:} say the Marwadi for a lesson, then the Marwadi for to learn by heart, then say \emph{uchārṇo} again
\end{itemize}

\begin{tcolorbox}[breakable,colback=teal!4,colframe=teal!35!black,title={Before you move on}]
What does \mw{उचारणो} mean? (\textquotedblleft{}To pronounce\textquotedblright{}.) Say it once more.
\end{tcolorbox}
